\documentclass[a4paper,12pt]{extarticle}
\usepackage{mathptmx}
\usepackage{cite}
\usepackage{amsmath,amssymb,amsfonts}
\usepackage{graphicx}
\usepackage{textcomp}
\usepackage{booktabs}
\usepackage{float}
\usepackage{makecell}
\usepackage[colorlinks=true, allcolors=black]{hyperref}
\usepackage{subcaption}
\usepackage[table]{xcolor}
\usepackage{multirow}
\usepackage{tabularx}
\usepackage{array}
\usepackage{setspace}
\usepackage[top=1.2in, bottom=1in, left=1.2in, right=1.0in]{geometry}
\usepackage{placeins}   % for \FloatBarrier

% --- Float placement tuning: force top/bottom only ---
\setcounter{topnumber}{2}
\setcounter{bottomnumber}{2}
\setcounter{totalnumber}{4}
\renewcommand{\topfraction}{0.9}      % allow up to 90% of page top for floats
\renewcommand{\bottomfraction}{0.9}   % allow up to 90% of page bottom for floats
\renewcommand{\textfraction}{0.1}     % require only 10% of page for text
\renewcommand{\floatpagefraction}{0.5}% discourage float-only pages
% --------------------------------------------------------

\setstretch{1.2}
\setlength{\parskip}{0.6em}
\setlength{\parindent}{0em}
\setlength{\arrayrulewidth}{0.4mm}
\setlength{\tabcolsep}{6pt}
\renewcommand{\arraystretch}{1.3}

\raggedbottom   % flexible page height, helps floats find a place

\title{\textbf{3D Mosque}}
\author{}
\date{}

\begin{document}

\begin{titlepage}
% ... (title page unchanged) ...
\begin{center}

{\fontsize{14pt}{16.8pt}\selectfont\textbf{CSE-4208 : Computer Graphics Laboratory}}\\[0.9\baselineskip]

{\fontsize{16pt}{18pt}\selectfont\textbf{3D Mosque}}\\[0.5\baselineskip]

{\fontsize{11pt}{13pt}\selectfont A 3D Simulation using Modern Opengl}\\[1cm]

\includegraphics[height=1.2in]{logo}\\[1cm]

{\large
\renewcommand{\arraystretch}{1.1}
\hspace*{0.9cm}

\begin{tabularx}{0.9\linewidth}{p{0.42\linewidth} @{
\hspace{1.5cm}} p{0.42\linewidth}}

\textbf{Submitted to:} & \textbf{Submitted by:} \\

Dr. Sk. Md. Masudul Ahsan & Shoab Ahamed \\

Professor & Roll: 2007015 \\

Department of CSE, KUET &  \\

\\

Md Tajmilur Rahman  \\

Lecturer  \\

Department of CSE, KUET &  \\

\end{tabularx}
}

\vfill

\textbf{Department of Computer Science and Engineering}\\
\textbf{Khulna University of Engineering \& Technology}\\
\textbf{Khulna 9203, Bangladesh}\\
\textbf{April 12, 2026}

\end{center}
\end{titlepage}

\clearpage
\pagenumbering{arabic}

\section{Objectives}
% ... (unchanged text) ...
This project focuses on building an interactive, real-time 3D mosque scene with concise, demonstrable goals:
\begin{itemize}
    \item Create a modular, extendable 3D mosque scene using modern OpenGL in C++ with reusable primitive meshes.
    \item Implement controllable lighting and shading (Phong), and basic material/texture blending for visual clarity.
    \item Provide responsive camera controls and simple interactive animations (doors, fans, water) for a demonstrable user experience.
\end{itemize}

\section{Introduction}
% ... (unchanged text) ...
This project presents an interactive 3D mosque simulation developed with modern OpenGL and C++. The goal is to build a complete architectural scene that is visually clear, modular in design, and easy to extend during future development. The implementation combines reusable primitive geometry, custom mesh generation, texture mapping, and real-time user interaction.

The rendering pipeline is built using GLFW for window and input management, GLAD for OpenGL function loading, and GLM for matrix/vector operations. Scene construction follows a component-based approach where major parts (mosque body, mihrab, minarets, doors, bookshelves, chandelier, wadhu area, and water effects) are designed as reusable drawing units.

From a user perspective, the application supports camera movement, light toggling, shading mode switching, and animated elements such as doors, windows, fans, and faucet water flow. This makes the scene useful not only as a visual model but also as a practical demonstration of core real-time graphics concepts taught in the course.

\section{Methodology}
% ... (unchanged text) ...
This project uses practical real-time rendering techniques without overly complex theory. The focus is on clear implementation choices that directly support scene quality and interaction.

\subsection*{Primitive-Based Modeling}
Most scene objects are built by combining reusable primitive meshes: cube, cylinder, sphere/hemisphere, and cone. Cubes are used for walls, shelves, doors, windows, and road blocks; cylinders are used for pillars, minaret shafts, chandelier rods, and faucets; spheres are used for domes and decorative bulbs; cones are used for finials and top ornaments. By scaling, rotating, and translating these primitives hierarchically, complex architectural structures are created without requiring heavy external models.

\subsection*{Phong Shading and Light Types}
The scene uses a Phong-based shader pipeline as the primary lighting model. Light contribution is controlled through three light types to make the scene visually balanced and easy to analyze:
\begin{itemize}
    \item \textbf{Directional light:} Represents a global sun-like source affecting the entire scene uniformly by direction.
    \item \textbf{Point lights:} Used for local lamps and decorative light sources; intensity decreases with distance.
    \item \textbf{Spot light:} Used as a focused hanging light source inside the mosque with cone-based direction.
\end{itemize}

\subsection*{Bezier Curve for Mihrab Construction}
The mihrab profile is created from a cubic Bezier curve and then extruded to form a smooth curved surface. The curve sampling equation used is:
$$
\mathbf{B}(t) = (1-t)^3\mathbf{P}_0 + 3(1-t)^2 t\mathbf{P}_1 + 3(1-t)t^2\mathbf{P}_2 + t^3\mathbf{P}_3, \quad 0 \le t \le 1.
$$
Sampled curve points are connected into a mesh, and then wall/roof/floor faces are generated for the final mihrab geometry.

\subsection*{Image Loading and Texture Setup}
Textures are loaded using \texttt{stb\_image}. For each image, the program reads width, height, and channel count, chooses the OpenGL format (RGB/RGBA), uploads data with \texttt{glTexImage2D}, and generates mipmaps with \texttt{glGenerateMipmap}. Wrapping and filtering modes are configured during loading so each object can reuse the same texture-loading function with different visual behavior.

\section{Implementation}
% ... (unchanged text) ...
This section summarizes the major objects and features, how they were built from primitive geometry, and where screenshots can be inserted.

\begin{figure}[t]   % explicitly force TOP placement
\centering
\includegraphics[width=\linewidth, height=0.75\textheight, keepaspectratio]{mosque_full1.png}
\caption{Full project scene}
\label{fig:full_scene}
\end{figure}

Key object groups included in the scene are:
\begin{itemize}
    \item Main mosque body
    \item Mihrab (Bezier geometry)
    \item Dome
    \item Minarets
    \item Doors and windows
    \item Chandelier
    \item Corner bookshelves
    \item Back-wall alcove bookshelves
    \item Trees
    \item Roads and ponds
    \item Wadhu faucets and water droplets
    \item Ceiling fans
    \item Directional, point, and spot lights
\end{itemize}

\subsection*{Mihrab (Bezier Extrusion)}
The mihrab is generated from a Bezier profile and extruded to create a smooth curved niche. The control points are arranged to shape the arch opening and side curvature, then sampled profile points are connected into triangles along the extrusion depth. Cap surfaces and UV mapping are added so the front face, inner curve, and side walls all receive consistent texturing.
\begin{figure}[t]   % top placement
\centering
\includegraphics[width=0.65\linewidth, height=0.7\textheight, keepaspectratio]{mihrab.png}
\caption{Mihrab}
\label{fig:impl_mihrab}
\end{figure}

\subsection*{Dome}
The dome is built from hemisphere-based sphere geometry with dome-specific scaling and texture mapping for a clean top silhouette. After scaling, it is positioned as an upper shell above the main prayer hall so its base aligns with the roof ring. Material settings are tuned to keep highlights soft, which helps preserve the rounded form under multiple lights.
\begin{figure}[t]   % top placement
\centering
\includegraphics[width=0.65\linewidth, height=0.7\textheight, keepaspectratio]{dome.png}
\caption{Dome}
\label{fig:impl_dome}
\end{figure}

\subsection*{Minaret}
Each minaret is assembled from stacked cubes/cylinders plus dome and cone ornaments. Hierarchical transforms are used to place repeated structural parts, so one local design can be reused at multiple corners with different world transforms. This keeps proportions consistent while reducing duplicated modeling code.
\begin{figure}[t]   % top placement
\centering
\includegraphics[width=0.65\linewidth, height=0.7\textheight, keepaspectratio]{minarets.png}
\caption{Minaret}
\label{fig:impl_minaret}
\end{figure}

\subsection*{Door}
Doors are built from cube panels with hinge-based rotation. Separate left/right leaves allow open-close animation controlled at runtime. Frame members and panel offsets are modeled independently, so each leaf rotates around an edge pivot without breaking alignment with the doorway.
\begin{figure}[t]   % top placement
\centering
\includegraphics[width=0.65\linewidth, height=0.7\textheight, keepaspectratio]{door_open.png}
\caption{Door }
\label{fig:impl_door}
\end{figure}

\subsection*{Window}
Windows are modeled similarly to door leaves using cube-based panels and frame segments. Each window leaf rotates around its hinge axis to create a realistic open-close motion and to support runtime interaction. Glass and frame parts use different material settings so openings remain visually readable from both interior and exterior views.
\begin{figure}[t]   % top placement
\centering
\includegraphics[width=0.65\linewidth, height=0.7\textheight, keepaspectratio]{window_open.png}
\caption{Window }
\label{fig:impl_window}
\end{figure}

\subsection*{Fan}
Ceiling fans are assembled from a cylinder downrod, a central hub, and three blade meshes (cube primitives). The fan rotates by continuously increasing a blade angle value over time ($\theta = \theta + \omega\Delta t$), then applying this angle as a Y-axis rotation transform each frame. Blade meshes are placed with equal angular offsets around the hub, so one update variable drives synchronized motion for the full assembly.
\begin{figure}[t]   % top placement
\centering
\includegraphics[width=0.65\linewidth, height=0.7\textheight, keepaspectratio]{fans.png}
\caption{Ceiling fan }
\label{fig:impl_fan}
\end{figure}

\subsection*{Animation Logic (Door, Window, Fan)}
Door and window states are toggled from keyboard input and updated smoothly using interpolation between closed and open targets each frame, which avoids abrupt jumps in motion. Fan motion is time-driven and independent of camera movement, so rotation speed remains visually consistent across frame rates. The update path is centralized in per-frame logic, keeping interaction state changes separate from drawing calls and making additional animated objects easier to integrate.

\subsection*{Chandelier}
The chandelier is made from cylinders/spheres with a central hanging structure and circular ring lamps. It contributes both visual detail and local illumination. Ring segments and hanging rods are composed with local transforms, then mirrored around the center to maintain a balanced shape from multiple camera angles.
\begin{figure}[t]   % top placement
\centering
\includegraphics[width=0.65\linewidth, height=0.7\textheight, keepaspectratio]{chandelier.png}
\caption{Chandelier }
\label{fig:impl_chandelier}
\end{figure}

\subsection*{Bookshelves (Two Types)}
Two bookshelf designs are used: low-height corner bookshelves and recessed back-wall alcove bookshelves. Both are made from cube-based panel structures with textured book blocks. The corner variant prioritizes compact footprint near wall edges, while the alcove variant fills recessed wall space with layered shelves and denser book placement.
\begin{figure}[t]   % top placement
\centering
\begin{subfigure}{0.48\textwidth}
\centering
\includegraphics[width=\linewidth, height=0.55\textheight, keepaspectratio]{open_bookshelf.png}
\caption{Corner bookshelf}
\end{subfigure}
\hfill
\begin{subfigure}{0.48\textwidth}
\centering
\includegraphics[width=\linewidth, height=0.55\textheight, keepaspectratio]{aclove_bookshelf.png}
\caption{Back-wall alcove bookshelf}
\end{subfigure}
\caption{Bookshelf types (replace with implementation screenshots)}
\label{fig:impl_bookshelves}
\end{figure}

\subsection*{Tree}
Trees are created using a simple procedural composition of primitive geometry, where trunk and branch structures are formed with scaled cylinders and foliage is approximated with clustered geometry. This approach keeps rendering lightweight while adding environmental depth around the mosque. Branch segments are scaled and rotated with small variation to avoid a repetitive silhouette while still using low-cost meshes.
\begin{figure}[t]   % top placement
\centering
\includegraphics[width=0.65\linewidth, height=0.7\textheight, keepaspectratio]{trees.png}
\caption{Tree object }
\label{fig:impl_tree}
\end{figure}

\subsection*{Road and Pond}
The outer environment includes road segments and pond surfaces to provide spatial context and realism. Roads are constructed from large scaled primitives with repeating texture coordinates, while pond surfaces use water-like texturing and material settings to distinguish them from surrounding ground. Both elements are aligned with terrain-level transforms so they integrate naturally with nearby trees, mosque boundaries, and walkway edges.
\begin{figure}[t]   % top placement
\centering
\includegraphics[width=0.75\linewidth, height=0.7\textheight, keepaspectratio]{pond.png}
\caption{Road and pond view }
\label{fig:pond}
\end{figure}

\subsection*{Faucet and Water (Side by Side)}
Faucets are modeled mainly with cylinders and small cube details. Water is implemented as a particle-like droplet system emitted from faucet source points. Briefly, each droplet stores position, velocity, lifetime, and source index; every frame gravity updates velocity, position is advanced over $\Delta t$, and droplets are respawned when lifetime expires or when they reach the channel zone, producing continuous flowing water. Because droplets are recycled from a fixed pool, the effect remains stable in real time while still giving a continuous flow appearance beside the faucet geometry.
\begin{figure}[t]   % top placement
\centering
\begin{subfigure}{0.48\textwidth}
\centering
\includegraphics[width=\linewidth, height=0.55\textheight, keepaspectratio]{faucet_only.png}
\caption{Faucet model}
\end{subfigure}
\hfill
\begin{subfigure}{0.48\textwidth}
\centering
\includegraphics[width=\linewidth, height=0.55\textheight, keepaspectratio]{faucet_only.png}
\caption{Water particle flow}
\end{subfigure}
\caption{Faucet and water visualization (replace with implementation screenshots)}
\label{fig:impl_faucet_water}
\end{figure}

\section{Discussion}
% ... (unchanged text) ...
The implementation emphasizes clarity and modularity: primitive mesh classes, a simple shader manager, and a clear separation between scene assembly and rendering code. This makes the project easy to extend and to reuse geometry for other scenes.

Limitations:
\begin{itemize}
    \item No physically-based rendering or shadow mapping implemented; lighting is analytic and less physically accurate than PBR pipelines.
    \item Performance optimizations (frustum culling, geometry instancing for large numbers of props, LODs) are minimal and could be improved for larger scenes.
    \item Water reflections and advanced post-processing are not included in this iteration.
\end{itemize}

Possible future work includes adding shadow mapping, PBR materials, GPU particle systems, and scene-level culling/instancing to improve both fidelity and performance.

\section{Conclusion}
% ... (unchanged text) ...
This project presents a modular OpenGL-based simulation of a mosque complex demonstrating geometry generation, shader-based lighting, texture mapping, and simple interactive animations. It serves as a solid educational platform for exploring real-time rendering concepts and can be extended with modern effects such as shadows and PBR in future work.

\FloatBarrier   % force all remaining floats to be placed here, before references

\section{References}
\begin{itemize}
    \item LearnOpenGL — Jo\'s de Vries. Available: \url{https://learnopengl.com}
    \item OpenGL Architecture Review Board, \textit{OpenGL 3.3 Core Profile Specification}.
    \item G. Sellers, R. Wright, and N. Haemel, \textit{OpenGL SuperBible}, 7th ed., Addison-Wesley, 2015.
    \item M. A. P. S. M. Kilgard, \textit{OpenGL Programming Guide} ("The Red Book").
    \item GLFW Documentation. Available: \url{https://www.glfw.org/docs/latest/}
    \item GLM (OpenGL Mathematics). Available: \url{https://github.com/g-truc/glm}
    \item stb\_image — Sean Barrett. Available: \url{https://github.com/nothings/stb}
    \item A. K. Peters, \textit{Real-Time Rendering} (for foundational concepts and best practices).
\end{itemize}

\end{document}